The scale of the mixed derivative is the obstruction to a direct
absolute-error argument.  With $a_b=S_b''(0)>0$ and
$q_p=\sqrt{a_0a_p}$, the quadratic reference is
\begin{equation}\label{eq:v50-reference-scale}
 d_j^0=-W_{j,uv}(0,0)=\frac{q_p}{\sinh(j\gamma)}.
\end{equation}
An absolute remainder $O(\theta^j)$ would give only
$O((\theta e^\gamma)^j)$ after division by this reference; its
decay is not automatic.  Instead the proof uses the exact cofactor
identity
\begin{equation}\label{eq:v50-relative-log}
 \log b_j=\sum_{i=0}^{j-1}
 \log\{g[-\ell_{i\bmod2,uv}(y_i,y_{i+1})]\}
                   -\log\det(I+G_j\Delta H_j).
\end{equation}
Here $y$ is the stationary bridge, $G_j$ is the inverse reference
interior Hessian, and $\Delta H_j$ is the subtracted nonlinear
Hessian.  The trace-norm estimates sum the localized endpoint layers,
whose weights have bounded sum independent of $j$, rather than
multiplying an unweighted entry bound by the number of collisions.
Two separated compressions of the reference Green operator then
converge to the half-line operators.  The trace-series comparison in
\eqref{eq:v4-det-tail}--\eqref{eq:v4-limit-integral} yields a
relative amplitude estimate before the exponentially small scale is
removed.  The long-flight limit is taken at fixed positive excess
$d$, not simultaneously with a zero-offset quadratic approximation.
This is why it retains nonlinear actions rather than only curvature.

